\section{Anchored prefix fibres for proper-face CRT flags}
\label{sec:anchored-prefix-flagged-crt}

The residue-cube formulation in
Section~\ref{sec:flagged-crt-residue-cube} turns an FCRT flag into a nested
cross-fibre problem.  This section extracts a finite positive result from
Bae's anchored-fibre selection lemma~\cite{Bae2026Erdos684}.  An ordered
reservoir with sufficiently many prefix samples, relative to its residue
subgroup and weight-cell entropy, contains a light zero-sum interval.  Its
complement is the required proper face.  The entropy hypothesis is explicit;
no uniform endpoint theorem supplying it is assumed.

Let $p^{e_p}\Vert c$, with $e_p\ge2$, be a target source and let $S$ be a
source face disjoint from $p$.  For a sink packet $A\subseteq J$, retain the
notation
\[
 G_p=(\mathbb Z/p^{e_p}\mathbb Z)^\times,
 \qquad
 \Phi_p(A)=\prod_{q\in A}g_{p,q},
 \qquad
 y(A)=\sum_{q\in A}\log q.
\tag{APF-1}\label{eq:apf-notation}
\]
Here $g_{p,q}$ contains the complete $q$-power on the $a$-arm and its
inverse on the $b$-arm.  Thus
\[
 \Phi_p(J)=-1,
 \qquad
 p^{e_p}\mid a_A+b_A
 \quad\Longleftrightarrow\quad
 \Phi_p(A)=-1.
\tag{APF-2}\label{eq:apf-compatibility}
\]
Fix a packet $T\subseteq J$ satisfying the joint conditions
\[
 \Phi_S(T)=-1,\qquad \Phi_p(T)=-1,\qquad x(S)\le y(T).
\tag{APF-3}\label{eq:apf-joint-block}
\]
The second condition is stronger than ordinary FCRT block compatibility,
but it is automatic when $T=J$.

\subsection{The complement dictionary}

Call a nonempty $D\subseteq T$ a $p$-zero-sum deletion if
$\Phi_p(D)=1$.

\begin{proposition}[Zero-sum deletion and compatible proper face]
\label{prop:apf-complement-dictionary}
The assignments
\[
                         D\longmapsto U=T\setminus D,
 \qquad U\longmapsto D=T\setminus U
\tag{APF-4}\label{eq:apf-complement-map}
\]
give inverse correspondences between nonempty $p$-zero-sum deletions of
$T$ and nonempty proper subfaces $U\subsetneq T$ satisfying
$\Phi_p(U)=-1$.
\end{proposition}

\begin{proof}
If $D$ and $U=T\setminus D$ partition $T$, then
\[
                         \Phi_p(T)=\Phi_p(U)\Phi_p(D).
\]
Thus $\Phi_p(T)=-1$ and $\Phi_p(D)=1$ imply $\Phi_p(U)=-1$.
Nonemptiness of $D$ makes $U$ proper.  If $U$ were empty, then $D=T$ and
$\Phi_p(T)=1$, contradicting $\Phi_p(T)=-1$.  The two group elements are
distinct because $p^{e_p}\ge4$.

Conversely, if $U\subsetneq T$ and both $U$ and $T$ have target label $-1$,
then cancellation in $G_p$ gives $\Phi_p(T\setminus U)=1$.  Properness makes
$T\setminus U$ nonempty, and the two complement operations in
\eqref{eq:apf-complement-map} are inverse.
\end{proof}

The block surplus is $\sigma=y(T)-x(S)$.  If $U=T\setminus D$, its capped
FCRT credit is
\[
 \rho=\min\{\sigma,y(U)\}
     =\min\{y(T)-x(S),\,y(T)-y(D)\}.
\tag{APF-5}\label{eq:apf-deletion-credit}
\]
It follows directly that the full surplus is reusable precisely when
\[
                              y(D)\le x(S).
\tag{APF-6}\label{eq:apf-full-surplus-condition}
\]

\subsection{Anchored weighted selection}

Choose a nonempty reservoir $K\subseteq T$, order it as
$K=(q_1,\ldots,q_m)$, and put
\[
 H=\langle g_{p,q}:q\in K\rangle\le G_p,
 \qquad W=y(K).
\tag{APF-7}\label{eq:apf-reservoir}
\]
Then $m\ge1$ and $W>0$.  Fix $L>0$, set
\[
                         N=\left\lceil\frac WL\right\rceil,
\tag{APF-8}\label{eq:apf-cell-count}
\]
and partition $[0,W]$ into $N$ consecutive cells of diameter at most $L$.
Let $\mathcal F\subseteq\{1,\ldots,m\}$ be a set of forbidden prefix gaps.

\begin{theorem}[Anchored prefix flag]
\label{thm:apf-anchored-prefix}
Under \eqref{eq:apf-joint-block}, assume the exact entropy inequality
\[
                  m+1>|H|N\bigl(|\mathcal F|+1\bigr).
\tag{APF-9}\label{eq:apf-entropy}
\]
There are $0\le i<j\le m$, with $j-i\notin\mathcal F$, for which
\[
 D=\{q_{i+1},\ldots,q_j\}\ne\varnothing,\qquad
 \Phi_p(D)=1,\qquad y(D)\le L.
\tag{APF-10}\label{eq:apf-light-deletion}
\]
Consequently $U=T\setminus D$ is a legal nonempty proper target face and
\[
 \rho\ge
 \min\{y(T)-x(S),\,y(T)-L\}.
\tag{APF-11}\label{eq:apf-credit-lower-bound}
\]
If $L\le x(S)$, then $\rho=y(T)-x(S)$, so the complete block surplus is
reused.
\end{theorem}

\begin{proof}
For $0\le k\le m$, define the prefix packet, residue, and weight by
\[
 P_k=\{q_1,\ldots,q_k\},
 \qquad R_k=\Phi_p(P_k),
 \qquad Z_k=y(P_k).
\]
Code $k$ by the pair consisting of $R_k\in H$ and the weight cell containing
$Z_k\in[0,W]$.  There are at most $|H|N$ codes.  Bae's anchored-fibre lemma
states that a map on $\{0,\ldots,m\}$ with at most $C$ codes has equal-code
indices $i<j$ whose positive difference avoids $\mathcal F$, provided
$m+1>C(|\mathcal F|+1)$~\cite[Lemma~5.1]{Bae2026Erdos684}.  Apply it with
$C=|H|N$.

Equality of residue coordinates gives
\[
 1=R_jR_i^{-1}
   =\prod_{r=i+1}^{j}g_{p,q_r}=\Phi_p(D).
\]
Equality of weight cells and monotonicity of the prefix weights give
$y(D)=Z_j-Z_i\le L$.  Since $i<j$, the deletion is nonempty.
Proposition~\ref{prop:apf-complement-dictionary} turns it into a legal face.
Moreover,
$y(U)=y(T)-y(D)\ge y(T)-L$, and
\eqref{eq:apf-deletion-credit} proves
\eqref{eq:apf-credit-lower-bound}.  When $L\le x(S)$,
$y(T)-L\ge y(T)-x(S)$, so the first entry is the minimum and the full
surplus is reused.
\end{proof}

The forbidden-gap clause is retained exactly.  It may encode a bad tail,
unwanted interval size, or another obstruction depending only on $j-i$.
The theorem does not assert that endpoint arithmetic supplies a uniformly
small forbidden set.

\begin{corollary}[Light-reservoir full-surplus criterion]
\label{cor:apf-light-reservoir}
If
\[
                         |K|\ge|H|,\qquad y(K)\le x(S),
\tag{APF-12}\label{eq:apf-light-reservoir}
\]
then the complete surplus of $(S,T)$ can be sent to $p$.
\end{corollary}

\begin{proof}
Take $L=x(S)$ and $\mathcal F=\varnothing$.  The hypotheses imply
$0<W\le L$, so $N=1$, and $|K|+1>|H|$.
Theorem~\ref{thm:apf-anchored-prefix} applies.
\end{proof}

\begin{corollary}[Full packet and the scalar boundary]
\label{cor:apf-full-packet}
Suppose $I\setminus\{p\}=S\ne\varnothing$, $x(S)\le Y=y(J)$, and the
hypotheses of Theorem~\ref{thm:apf-anchored-prefix} hold with $T=J$ and
$L\le x(S)$.  The resulting one-block configuration has
\[
 B_{\rm FCRT}
  =\bigl(x_p-(Y-x(S))\bigr)_+
  =(X-Y)_+.
\tag{APF-13}\label{eq:apf-scalar-boundary}
\]
\end{corollary}

\begin{proof}
The full packet has label $-1$ in every source coordinate.  The selected
block covers $S$, consumes all sinks, and sends its full surplus
$Y-x(S)$ to the sole residual source $p$.  Its clipped residual is the first
quantity in \eqref{eq:apf-scalar-boundary}, while
\[
                  x_p-(Y-x(S))=x_p+x(S)-Y=X-Y.
\]
The universal once-charge lower bound proves optimality.
\end{proof}

This corollary removes a finite fragmentation loss when its entropy premise
holds.  It does not estimate $(X-Y)_+$ uniformly and therefore is not an
$abc$ theorem.

\subsection{Recovery of the two motivating flags}

For $(a,b,c)=(1,675,676)$, choose
\[
 S=\{13\},\qquad p=2,\qquad T=J=\{3,5\},\qquad K=\{5\}.
\]
Modulo $4$, the $b$-arm generator satisfies $g_{2,5}=5^{-2}=1$.
Consequently $H$ is trivial, $|K|=|H|=1$, and
$y(K)=\log5<\log13=x(S)$.  Corollary~\ref{cor:apf-light-reservoir}
produces $D=\{5\}$ and $U=\{3\}$, recovering the full-surplus flag and the
boundary factor $26/15$.

For $(1,65024,65025)$, choose
\[
 S=\{5,17\},\qquad p=3,\qquad
 T=J=\{2,127\},\qquad K=\{127\}.
\]
Modulo $9$, $g_{3,127}=127^{-1}=1$.  With
$L=\log127$ and $\mathcal F=\varnothing$, one has $m=N=|H|=1$, so the
strict entropy condition is $2>1$.  Theorem~\ref{thm:apf-anchored-prefix}
gives $D=\{127\}$ and $U=\{2\}$.  Therefore
\[
 \rho\ge
 \min\left\{\log\frac{254}{85},\log\frac{254}{127}\right\}
 =\log2.
\]
The face itself has weight $\log2$, so equality holds.  This is the
cap-active flag giving FCRT boundary factor $3/2$.

\subsection{A complete-premise failure of the naive Boolean shortcut}

Consider the assertion
\[
 \tag{NBF}\label{eq:apf-nbf}
 2^{|J|}>|G_p|
 \quad\Longrightarrow\quad
 \text{there is a nonempty proper $U\subsetneq J$ with $\Phi_p(U)=-1$}.
\]
Its proposed FCRT consequence would give a target-$p$ flag for a saturated
full block whose source face is disjoint from $p$.

\begin{proposition}[Complete-premise NBF counterexample]
\label{prop:apf-nbf-counterexample}
The primitive endpoint
\[
                         (a,b,c)=(1,4715,4716)
\]
satisfies the Boolean inequality and has a compatible saturated full block
disjoint from the target $p=3$, but it has no nonempty proper
$3$-compatible packet.  Hence \textup{NBF} is false.
\end{proposition}

\begin{proof}
The exact factorizations are
\[
 4715=5\cdot23\cdot41,
 \qquad
 4716=2^2\cdot3^2\cdot131.
\tag{APF-14}\label{eq:apf-4715-factorization}
\]
Thus $p^2=9\mid c$, $J=\{5,23,41\}$, and
\[
             2^{|J|}=8>6=|(\mathbb Z/9\mathbb Z)^\times|=|G_p|.
\]
The full block $S=\{2\}$, $T=J$ is disjoint from $p$, compatible because
$4\mid4716$, and saturated because $x(S)=\log2<\log4715=y(J)$.

Each sink prime is congruent to $5$ modulo $9$ and occurs to the first power
on the $b$-arm.  Therefore
\[
                         g_{3,q}=q^{-1}=2\pmod9
                         \qquad(q\in J),
\]
and every packet has label
\[
                         \Phi_3(U)=2^{|U|}\pmod9.
\]
A nonempty proper packet has size one or two, hence label $2$ or $4$, never
$-1=8$.  Equivalently, its arm sum is congruent to $6$ or $8$, rather than
zero, modulo $9$.  Only the full packet, of size three, is compatible.
Every possible FCRT face is a nonempty proper subset of $J$, so no flag can
target $3$.
\end{proof}

The Boolean collision predicted by pigeonhole still occurs: all singleton
packets have label $2$, and all doubletons have label $4$.  Distinct packets
of the same size are incomparable, so cancellation produces a signed
relation rather than a positive zero-sum deletion.  The counterexample does
not meet Theorem~\ref{thm:apf-anchored-prefix}'s entropy premise.  Every
nonempty reservoir generates $H=\langle2\rangle=G_p$, so $|H|=6$ while
$|K|\le3$ and $N\ge1$.

Accordingly, this example retires only the implication \textup{NBF}.  It
does not refute the anchored theorem, FCRT-1, SCRT-0, the residue-cube route,
or the standard $abc$ conjecture.

\subsection{The remaining uniform bottleneck}

For a jointly compatible packet $T$, define the weighted zero-sum radius
\[
 \zeta_p(T)=
 \min\{y(D):\varnothing\ne D\subseteq T,\ \Phi_p(D)=1\},
\tag{APF-15}\label{eq:apf-zero-sum-radius}
\]
with value $+\infty$ when no such deletion exists.  The complement dictionary
identifies proper target faces with these deletions, and full surplus reuse
is equivalent to $\zeta_p(T)\le x(S)$.

The anchored theorem controls this radius only after an ordered reservoir
satisfies
\[
 |K|+1>
 |\langle g_{p,q}:q\in K\rangle|
 \left\lceil\frac{y(K)}L\right\rceil
 (|\mathcal F|+1),
\tag{APF-16}\label{eq:apf-uniform-bottleneck}
\]
with $L$ small enough to yield useful credit.  No uniform endpoint theorem
forcing \eqref{eq:apf-uniform-bottleneck} is known.  The generated subgroup
may have prime-power scale, whereas $|K|$ counts only distinct sink primes.
Bae's application has an exponentially longer sample interval and a
separate exponential tail estimate~\cite{Bae2026Erdos684}; endpoint analogues
of those inputs have not been proved.  The NBF counterexample shows why
replacing the ordered prefixes by all Boolean packets does not by itself
solve the mismatch.

The remaining positive problem is a weighted zero-sum or Davenport-type
estimate using the special endpoint relation $\Phi_p(J)=-1$, or an arithmetic
theorem forcing many light sink generators into a small subgroup.  Absence
of such an estimate records the present bottleneck and does not retire the
parent route.

\subsection{Lean boundary}

The companion module
\texttt{ABCAnchoredPrefixFlaggedCRT20260904} has $28$ public declarations,
and its audit file contains exactly $28$ corresponding
\texttt{\#print axioms} queries.  Both files compile with warnings treated as
errors.  The audit union is exactly \texttt{propext},
\texttt{Classical.choice}, and \texttt{Quot.sound}; no \texttt{sorry},
\texttt{admit}, \texttt{unsafe}, or \texttt{native\_decide} escape occurs.
The formalization proves the finite
anchored code-fibre avoidance lemma, additive prefix cancellation, the
abstract weighted selection theorem with a supplied finite cell code,
both directions of the complement dictionary, the two capped-credit
inequalities, the scalar residual identity, and the complete finite
arithmetic assertions for $(1,4715,4716)$.

The formal weighted theorem assumes a cell code together with the proved
implication that equal cells have interval-weight difference at most $L$.
It does not formalize the elementary real-ceiling construction of the cells
in \eqref{eq:apf-cell-count}.  It also does not construct the concrete unit
group $G_p$ uniformly from an arbitrary endpoint, prove the entropy estimate
\eqref{eq:apf-uniform-bottleneck}, or inhabit the uniform FCRT-1 gate.  The
two motivating examples are ordinary exact recoveries; the Lean module's
new endpoint arithmetic is the complete-premise $4715$ counterexample.
Thus the formal development contains neither an unconditional proof nor a
disproof of $abc$.
